% Generated draft: requires mathematical and source review.
\documentclass[11pt]{article}
\usepackage[margin=25mm]{geometry}
\usepackage{fontspec}
\setmainfont{Latin Modern Roman}
\usepackage{amsmath,amssymb,mathtools,unicode-math}
\setmathfont{Latin Modern Math}
\usepackage{enumitem,needspace,xurl,hyperref}
\setlength{\parindent}{0pt}
\setlength{\parskip}{5pt}
\setlength{\emergencystretch}{3em}
\newcommand{\N}{\mathbb{N}}
\newcommand{\Z}{\mathbb{Z}}
\newcommand{\Q}{\mathbb{Q}}
\newcommand{\R}{\mathbb{R}}
\newcommand{\C}{\mathbb{C}}
\newcommand{\F}{\mathbb{F}}
\newcommand{\A}{\mathbb{A}}
\newcommand{\PP}{\mathbb{P}}
\newcommand{\E}{\mathbb{E}}
\newcommand{\eps}{\varepsilon}
\newcommand{\dd}{\,\mathrm{d}}
\begin{document}
\section*{UnsolvedMath ID: 30004432}
\textbf{Source problem code:} catalogue entry
\subsection{Definitions and mathematical statement}
The supplied problem concerns long-range percolation on the one-dimensional lattice \(\Z\), indexed by an exponent \(s\). The record does not specify the edge probabilities, the convention relating \(s\) to their decay, or the percolation parameter and its critical value. Consequently, the following formulation is conditional on identifying the intended model and critical law.

For a specified such model, let \(\mathbb P_{s,c}\) denote its intended critical percolation law and define the two-point connectivity function by
\[
\tau_{s,c}(x,y)=\mathbb P_{s,c}(x\leftrightarrow y),\qquad x,y\in\Z,
\]
where \(x\leftrightarrow y\) means that an open path connects \(x\) and \(y\), with \(x\leftrightarrow x\) always holding. The usual critical triangle diagram at the origin is
\[
\Delta_s=\sum_{x,y\in\Z}\tau_{s,c}(0,x)\tau_{s,c}(x,y)\tau_{s,c}(y,0)\in[0,+\infty].
\]
The critical triangle condition is \(\Delta_s<\infty\).

Under this interpretation, the source asks whether
\[
\Delta_s=+\infty\qquad\text{for every }s>\frac43,
\]
and whether \(\Delta_{4/3}\) is finite or infinite. No upper bound on \(s\) is supplied. The record also does not explicitly identify the percolation parameter at which the triangle condition is evaluated; evaluation at criticality is the conventional interpretation used here and requires confirmation. A fully specified mathematical statement therefore requires the missing model, parameter range, and critical-law conventions.
\subsection{Short English statement}
For one-dimensional long-range percolation, determine whether the critical triangle diagram diverges for every \(s>4/3\), and determine its finiteness at \(s=4/3\). This formulation uses the usual critical triangle condition; the supplied record leaves the precise percolation model and exponent convention unspecified.
\subsection{Sources}
\begin{itemize}
\item Local UnsolvedMath registry: supplied record with permanent ID \(30004432\), title \emph{Triangle Condition in Long-Range Percolation}, exact source problem code \texttt{catalogue entry}, in the Oberwolfach Reports Open Problems collection. No individual proposer is supplied; the proposed year is recorded as 2020.
\item Supplied source citation: \emph{Mini-Workshop: One-sided and Two-sided Stochastic Descriptions} (2021), Oberwolfach Reports, pp.\ 601--637, DOI \texttt{10.4171/owr/2020/11}. \url{https://doi.org/10.4171/owr/2020/11}. The supplied background describes Oberwolfach Reports as documenting mathematical workshops and their open-problem discussions.
\item Supplied literature source, labeled OWR-2020-11: \url{https://ems.press/content/serial-article-files/46847}.
\item Supplied literature source: Heydenreich, van der Hofstad and Sakai, \emph{Mean-field behavior for long- and finite-range models}. \url{https://doi.org/10.1007/s10955-008-9580-5}.
\item Literature assessment checked on \texttt{2026-08-22}: the supplied curation labels the problem open and states that the triangle condition is proved for \(1<s<4/3\), while the questions for \(s>4/3\) and \(s=4/3\) remain unresolved in the primary literature it checked. This is supplied dated curation derived from the OWR material, not an independently verified literature assessment. Its status and cited claims require independent verification.
\end{itemize}
\end{document}
